\chapter{Korelasyon: İlişkinin Yönü ve Büyüklüğü}
\label{ch:correlation-regression}
\index{korelasyon}

\begin{hedefler}
Bu bölümün sonunda okuyucu:
\begin{itemize}
  \item iki değişken arasındaki ilişkiyi saçılım grafiğinde yön, biçim, güç ve aykırı değer bakımından inceleyebilecek,
  \item kovaryans ile Pearson korelasyonunun mantığını açıklayabilecek,
  \item Pearson ve Spearman korelasyonları arasında veri yapısına uygun seçim yapabilecek,
  \item korelasyon katsayısı, güven aralığı, $p$-değeri ve örneklem hacmini birlikte yorumlayabilecek,
  \item doğrusal olmayan ilişki, aykırı değer, ranj daralması ve alt grup karışımının korelasyona etkisini tanıyabilecek,
  \item korelasyon ile açıklanan varyans oranını birbirinden ayırabilecek,
  \item gözlenen ilişkiyi nedensellik iddiasına dönüştürmeden alternatif açıklamalar üretebilecek,
  \item Python ve SPSS çıktılarından bilimsel bir korelasyon raporu oluşturabilecektir.
\end{itemize}
\end{hedefler}

\begin{hazirlik}
Uyku süresi ile akademik başarı ilişkili bulunduğunda bu sonuç uykunun başarıyı
artırdığını kanıtlar mı? Ters nedensellik, ortak neden veya seçim sürecine
dayanan en az iki alternatif açıklama üretin.
\end{hazirlik}

\section{PDR vakası: Uyku süresi ve akademik kaygı}

Bir psikolojik danışma birimi, öğrencilerin gecelik ortalama uyku süresi ile
akademik kaygı puanları arasındaki ilişkiyi incelemektedir. Yirmi öğrenciden
elde edilen veride uyku süresi 4,5 ile 8,5 saat, kaygı puanı 40 ile 75 arasında
değişmektedir.

Araştırmacının yalnız ``korelasyon anlamlı mı?'' sorusunu sorması yeterli
değildir. Şu sorular birlikte ele alınmalıdır:

\begin{enumerate}
  \item İlişkinin yönü ve yaklaşık biçimi nedir?
  \item İlişki doğrusal mı, yoksa başka bir örüntü mü göstermektedir?
  \item Tek bir gözlem katsayıyı aşırı etkiliyor mu?
  \item İlişkinin büyüklüğü ve belirsizliği nedir?
  \item Gözlenen ilişki hangi nedensel ve nedensel olmayan süreçlerle açıklanabilir?
\end{enumerate}

Bu örnekte Pearson korelasyonu yaklaşık $r=-0{,}585$ bulunacaktır. Negatif
işaret, daha uzun uyku süresinin örneklemde daha düşük kaygı puanıyla birlikte
görülme eğiliminde olduğunu gösterir. Bu sayı henüz ilişkinin doğrusal
biçimini, aykırı gözlemleri veya nedensel yönünü kanıtlamaz.

\begin{researchbox}
Uyku ve kaygı öz-bildirimle aynı anda ölçülmüşse ölçüm hatası, ortak yöntem
etkisi ve zaman sırası belirsizliği vardır. Öğrencinin daha az uyuduğu için mi
kaygılandığı, kaygılandığı için mi uyuyamadığı veya sınav yükünün her ikisini
de mi etkilediği korelasyondan tek başına anlaşılamaz.
\end{researchbox}

\section{İlk adım: Saçılım grafiği}
\index{saçılım grafiği}
\index{doğrusal ilişki}

İki nicel değişken arasındaki ilişki önce saçılım grafiğiyle incelenir. Her
nokta bir öğrenciyi, yatay ve dikey konumlar o öğrencinin iki ölçümünü temsil
eder. Grafik dört temel özelliği görünür kılar; yalnız özet katsayılarına
dayanmak belirgin örüntüleri saklayabilir \cite{anscombe1973}:

\begin{description}
  \item[Yön] Değerler birlikte artıyor mu, biri artarken diğeri azalıyor mu?
  \item[Biçim] İlişki doğrusal, eğrisel, tekdüze veya kümeli mi?
  \item[Güç] Noktalar temel örüntünün çevresinde ne kadar sıkı toplanıyor?
  \item[Olağandışı gözlem] Aykırı veya yüksek etkili noktalar var mı?
\end{description}

\begin{figure}[htbp]
\centering
\resizebox{0.98\linewidth}{!}{%
\begin{tikzpicture}[alt={Pozitif, negatif ve doğrusal olmayan ilişkileri karşılaştıran üç saçılım grafiği},x=0.62cm,y=0.52cm]
  \foreach \dx/\title in {0/{pozitif},7/{negatif},14/{doğrusal olmayan}} {
    \draw[->,PrismGray] (\dx,0)--(\dx+5.2,0);
    \draw[->,PrismGray] (\dx,0)--(\dx,5.2);
    \node[font=\scriptsize,text=PrismGray] at (\dx+2.5,5.55) {\title};
  }
  \foreach \x/\y in {0.5/0.7,1/1.4,1.5/1.2,2/2.4,2.5/2.1,3/3.2,3.5/3.1,4/4.0,4.5/4.4}
    \fill[PrismBlue] (\x,\y) circle (2pt);
  \foreach \x/\y in {0.5/4.4,1/4.0,1.5/3.7,2/3.1,2.5/3.3,3/2.4,3.5/2.0,4/1.4,4.5/0.8}
    \fill[PrismSubheadingBlue] (7+\x,\y) circle (2pt);
  \foreach \x/\y in {0.4/4.2,0.9/3.1,1.4/2.0,1.9/1.1,2.4/0.7,2.9/0.9,3.4/1.5,3.9/2.6,4.4/4.1}
    \fill[PrismWarnOrange!85!black] (14+\x,\y) circle (2pt);
\end{tikzpicture}%
}
\caption{İlişkinin yönü ve biçimi saçılım grafiğinde görülür.}
\end{figure}

Son panelde güçlü bir U biçimli ilişki olmasına rağmen Pearson korelasyonu
sıfıra yakın olabilir. Bu nedenle ``$r=0$ ise değişkenler ilişkisizdir''
sonucu yanlıştır. Katsayı belirli bir ilişki türünü özetler; grafik daha geniş
örüntüyü gösterir.

\begin{mistakebox}
Saçılım grafiğine doğrusal eğilim çizgisi eklenmiş olması ilişkinin gerçekten
doğrusal olduğunu kanıtlamaz. Eğrilik, alt grup kümeleri ve uç gözlemler ayrıca
incelenmelidir.
\end{mistakebox}

\section{Kovaryanstan korelasyona}
\index{kovaryans}
\index{Pearson korelasyonu}

İki değişkenin ortalamalarından sapmaları çoğunlukla aynı işareti taşıyorsa
kovaryans pozitiftir; biri pozitifken diğeri negatif olma eğilimindeyse
kovaryans negatiftir:
\[
s_{xy}=\frac{1}{n-1}
\sum_{i=1}^n(x_i-\bar x)(y_i-\bar y).
\]

Kovaryansın büyüklüğü değişkenlerin ölçü birimlerine bağlıdır. Uyku saatten
dakikaya çevrildiğinde kovaryans 60 kat değişir. Pearson korelasyonu
kovaryansı iki standart sapmaya bölerek birimsiz hâle getirir:
\[
r=\frac{s_{xy}}{s_xs_y}
=\frac{\sum_{i=1}^n(x_i-\bar x)(y_i-\bar y)}
{\sqrt{\sum_{i=1}^n(x_i-\bar x)^2
             \sum_{i=1}^n(y_i-\bar y)^2}}.
\]

Standartlaştırılmış puanlarla eşdeğer gösterim
\[
r=\frac{1}{n-1}\sum_{i=1}^n z_{xi}z_{yi}
\]
biçimindedir. Aynı yönde standart sapma gösteren gözlemler pozitif, ters yönde
gösterenler negatif katkı yapar.

\section{Pearson korelasyonunun özellikleri}
\index{Pearson korelasyonu!özellikleri}

Pearson korelasyonu $-1$ ile $1$ arasındadır:

\begin{description}
  \item[$r=1$] Bütün noktalar pozitif eğimli bir doğru üzerindedir.
  \item[$r=-1$] Bütün noktalar negatif eğimli bir doğru üzerindedir.
  \item[$r=0$] Doğrusal ilişki yoktur; doğrusal olmayan ilişki olabilir.
\end{description}

Katsayının işareti yönü, mutlak değeri doğrusal ilişkinin derecesini gösterir.
X ile Y'nin yer değiştirmesi sonucu değiştirmez. Her iki değişkene sabit
eklemek veya pozitif bir sabitle çarpmak korelasyonu değiştirmez; değişkenlerden
birini negatif sabitle çarpmak yalnız işareti tersine çevirir.

Korelasyonun büyüklüğü eğim değildir. Aynı $r$ değeri farklı ölçü birimleri ve
farklı regresyon eğimleriyle birlikte görülebilir. Eğim, X'teki bir birimlik
değişime karşı Y'deki beklenen değişimi; korelasyon ise standartlaştırılmış
doğrusal birlikteliği özetler \cite{kutner2005,james2021}.

\begin{sezgibox}
Korelasyon, noktaların ince bir doğrusal şerit çevresinde ne kadar düzenli
dizildiğini anlatır. Şeridin dikliği ölçü birimine bağlıdır; korelasyonun
büyüklüğü ise noktaların şeride ne kadar sıkı bağlandığıyla ilgilidir.
\end{sezgibox}

\section{Adım adım PDR örneği}

Yirmi öğrencinin uyku ve kaygı verisinde Pearson katsayısı
\[
r=-0{,}585
\]
olarak hesaplanmıştır. Yorum üç aşamada yapılır.

\subsection*{1. Yön ve büyüklük}

İlişki negatiftir: daha uzun uyku süresi daha düşük kaygı puanıyla birlikte
görülme eğilimindedir. Mutlak değer yaklaşık 0,59'dur. Bu büyüklük PDR
bağlamında dikkate değer olabilir; ancak sabit ``orta'' veya ``güçlü'' etiketi
araştırma bağlamının yerine geçmemelidir.

\subsection*{2. Sıfır korelasyon testi}

$H_0:\rho=0$ için test istatistiği
\[
t=r\sqrt{\frac{n-2}{1-r^2}}
\]
olup $n-2$ serbestlik derecelidir. Örnekte
\[
t=-0{,}585\sqrt{\frac{18}{1-0{,}585^2}}
\approx-3{,}06,
\qquad p\approx0{,}007.
\]

Bu sonuç sıfır doğrusal korelasyon hipotezine karşı kanıt sunar. $p$-değeri
ilişkinin nedensel olduğunu veya gerçek katsayının tam $-0{,}585$ olduğunu
göstermez.

\subsection*{3. Güven aralığı}

Pearson korelasyonu için güven aralığı çoğunlukla Fisher dönüşümüyle kurulur:
\[
z_r=\frac12\ln\left(\frac{1+r}{1-r}\right),
\qquad
SE(z_r)=\frac{1}{\sqrt{n-3}}.
\]
Z ölçeğinde sınırlar oluşturulup ters dönüşüm uygulandığında yüzde 95 güven
aralığı yaklaşık
\[
[-0{,}816;\ -0{,}192]
\]
bulunur. Veri zayıf ile güçlü arasında değişebilen negatif evren
korelasyonlarıyla uyumludur; örneklemin küçük olması aralığı genişletmiştir.

\section{Korelasyon büyüklüğünü yorumlamak}
\index{korelasyon!yorumlama}

Korelasyon için evrensel küçük--orta--büyük sınırları yoktur. Yorum şu
özelliklere bağlıdır:

\begin{itemize}
  \item ölçümlerin güvenirliği ve ranjı,
  \item alan yazındaki tipik ilişkiler,
  \item ilişkinin hizmet veya karar açısından sonucu,
  \item örneklemin kimlerden oluştuğu,
  \item güven aralığının genişliği.
\end{itemize}

Çok küçük bir ilişki büyük bir öğrenci evreninde hizmet planlaması bakımından
önemli olabilir. Büyük görünen bir örneklem korelasyonu ise küçük ve seçilmiş
bir grupta kararsız olabilir. Kesin eşik yerine büyüklük ve belirsizlik birlikte
sunulmalıdır.

\section{Açıklanan varyans ve yanlış yorumlar}
\index{açıklanan varyans}
\index{R kare@$R^2$}

Basit doğrusal regresyonda tek bir yordayıcı ve kesme terimi bulunduğunda
\[
R^2=r^2
\]
ilişkisi vardır. Uyku--kaygı örneğinde
\[
r^2=(-0{,}585)^2\approx0{,}342.
\]
Bu değer, örneklemde kaygı puanlarındaki doğrusal değişkenliğin yaklaşık yüzde
34,2'sinin uyku süresiyle doğrusal olarak ortaklaştığını gösterir.

\begin{mistakebox}
``$r=-0{,}585$ olduğuna göre uyku kaygının yüzde 58,5'ini azaltır'' ifadesi
yanlıştır. Korelasyon yüzde değişim değildir. $r^2$ de tek başına nedensel
``açıklama'' oranı değildir; gözlemsel doğrusal ortak değişkenliği özetler.
\end{mistakebox}

Yüksek $R^2$ nedensellik, düşük $R^2$ de önemsizlik anlamına gelmez. Sonuç
değişkeninin doğal değişkenliği, ölçüm hatası ve karar bağlamı göz önünde
bulundurulmalıdır.

\section{Spearman sıra korelasyonu}
\index{Spearman korelasyonu}
\index{sıra korelasyonu}

Spearman korelasyonu ham puanların sıralarına Pearson korelasyonu uygulanarak
elde edilir. Bağlı sıra yoksa
\[
r_s=1-\frac{6\sum d_i^2}{n(n^2-1)}
\]
formülü kullanılabilir; uygulamada yazılım bağlı sıraları uygun biçimde ele
alır.

Spearman katsayısı şu durumlarda yararlıdır:

\begin{itemize}
  \item değişkenlerden biri anlamlı sıralara sahip fakat eşit aralıklı olmayan ölçekteyse,
  \item ilişki doğrusal değil ancak sürekli artan veya azalan, yani tekdüzeyse,
  \item uç değerler Pearson korelasyonunu aşırı etkiliyorsa,
  \item sıralama davranışı araştırma hedefiyle uyumluysa.
\end{itemize}

Uyku--kaygı verisinde Spearman korelasyonu yaklaşık
\[
r_s=-0{,}549,
\qquad p\approx0{,}012
\]
bulunur. Pearson ve Spearman sonuçlarının benzer olması negatif tekdüze
örüntünün yalnız birkaç uç değerden kaynaklanmadığına dair sınırlı bir
duyarlılık göstergesidir. İki katsayının farklı hedefleri olduğu unutulmamalıdır.

\begin{mistakebox}
Spearman korelasyonu ``varsayımsız Pearson'' değildir ve her aykırı değer
sorununu çözmez. Sıralar bilgi kaybına yol açabilir; çok sayıda bağlı sıra ve
olağandışı etkili gözlemler sonucu yine etkileyebilir.
\end{mistakebox}

\section{Doğrusal olmayan ve tekdüze olmayan ilişkiler}
\index{doğrusal olmayan ilişki}
\index{tekdüze ilişki}

Pearson doğrusal, Spearman ise tekdüze ilişkiyi özetler. X arttıkça Y önce
azalıyor sonra artıyorsa ilişki ne doğrusal ne de tekdüzedir. Her iki katsayı
sıfıra yakın olduğu hâlde güçlü bir U biçimli örüntü bulunabilir.

PDR bağlamında orta düzey stresin performansı artırıp çok düşük ve çok yüksek
stresin performansı düşürmesi böyle bir eğrisel ilişki oluşturabilir. Bu
durumda:

\begin{itemize}
  \item saçılım grafiği ve yumuşatılmış eğri incelenir,
  \item kuramsal olarak gerekçeli eğrisel terimler modellenir,
  \item doğrusal korelasyon ``ilişki yok'' diye yorumlanmaz.
\end{itemize}

İlişkinin biçimi veri görüldükten sonra sınırsız model denenerek seçilmemeli;
keşfedici bulgu sonraki çalışmada doğrulanmalıdır.

\section{Aykırı ve etkili gözlemler}
\index{aykırı değer}
\index{etkili gözlem}
\index{kaldıraç}

Pearson korelasyonu ortalama, standart sapma ve çapraz çarpımlara dayandığı
için aykırı gözlemlere duyarlıdır. X ekseninde diğer noktalardan çok uzakta
bulunan bir gözlem yüksek kaldıraçlıdır ve doğrusal örüntüyü güçlü biçimde
değiştirebilir.

Uyku--kaygı verisine 9 saat uyuduğunu ve kaygı puanının 90 olduğunu bildiren
tek bir öğrenci eklendiğinde Pearson korelasyonu yaklaşık $-0{,}585$'ten
$-0{,}153$'e düşmektedir. Spearman korelasyonu da yaklaşık $-0{,}338$'e
değişmektedir. Bu karşılaştırma gözlemin otomatik olarak silinmesini değil,
sonucun ona duyarlı olduğunu gösterir.

\begin{enumerate}
  \item Veri girişi, birim ve katılımcı kaydı doğrulanır.
  \item Gözlemin farklı bir alt grubu temsil edip etmediği incelenir.
  \item Grafik ve standartlaştırılmış değerler değerlendirilir.
  \item Gözlem dahil ve hariç sonuçlar duyarlılık analizi olarak sunulur.
  \item Silme kararı sonuçtan bağımsız, bilimsel gerekçeye dayandırılır.
\end{enumerate}

\begin{researchbox}
Uç bir kaygı puanı istatistiksel olarak etkili olduğu kadar mesleki açıdan da
önemli olabilir. Gerçek bir öğrenciyi yalnız korelasyonu ``bozduğu'' için
görünmez kılmak bilimsel ve etik açıdan savunulamaz.
\end{researchbox}

\section{Ranj daralması}
\index{ranj daralması}

Korelasyon örneklemde gözlenen değişkenlikten etkilenir. Yalnız yüksek başarı
puanlı öğrenciler seçildiğinde başarı ranjı daralır ve evrende var olan ilişki
örneklemde daha küçük görünebilir. Benzer biçimde yalnız yüksek risk grubunda
kaygı ile başka bir değişken arasındaki korelasyon zayıflayabilir.

Ranj daralması katsayının değişkenlerin değişmez bir özelliği olmadığını
gösterir. Örneklemin dahil edilme ölçütleri, minimum--maksimum değerleri ve
hedef evrenden nasıl ayrıldığı raporlanmalıdır.

\begin{sezgibox}
Bir ilişkinin görülebilmesi için her iki değişkende de yeterli çeşitlilik
gerekir. Yalnız aynı boya yakın kişileri seçerek boy ile kol uzunluğu
ilişkisini incelemek, cetvelin çok küçük bir bölümüne bakmaya benzer.
\end{sezgibox}

\section{Alt gruplar ve yanıltıcı toplam korelasyon}
\index{alt grup}
\index{toplam korelasyon}

İki veya daha fazla alt grup birlikte analiz edildiğinde toplam korelasyon,
grup içindeki ilişkilerden farklı hatta ters yönde olabilir. Örneğin lisans ve
lisansüstü öğrenciler hem yaş hem yardım arama puanı bakımından farklıysa,
grupların birleştirilmesi yaş ile yardım arama arasında yapay bir ilişki
üretebilir.

Bu durum bazen Simpson paradoksu olarak görünür. Araştırmacı:

\begin{itemize}
  \item saçılım grafiğinde noktaları anlamlı gruplara göre renklendirmeli,
  \item grup içi ve toplam ilişkileri karşılaştırmalı,
  \item alt grup analizini kuramsal gerekçeyle yürütmeli,
  \item çok sayıda veri sonrası alt grup aramanın yanlış pozitif riskini açıklamalıdır.
\end{itemize}

Alt grupları ayırmak her zaman doğru değildir; küçük grup örneklemleri çok
kararsız katsayılar üretebilir. Amaç, veri üretim yapısını görünür kılmaktır.

\section{Ölçüm güvenirliği ve korelasyon}
\index{ölçüm güvenirliği}
\index{korelasyon!zayıflama}

Rastgele ölçüm hatası gözlenen korelasyonu çoğunlukla sıfıra doğru zayıflatır.
İki yapının kuramsal ilişkisi güçlü olsa bile kısa, güvenirliği düşük ölçekler
daha küçük korelasyon verebilir. Bu nedenle kullanılan ölçme araçlarının
güvenirlik ve geçerlik kanıtları belirtilmelidir.

Yüksek güvenirlik ise korelasyonun geçerli veya nedensel olduğunu garanti
etmez. İki öz-bildirim ölçeği benzer soru dili ya da ortak yanıt eğilimi
nedeniyle yapay biçimde yüksek korelasyon gösterebilir.

\begin{mistakebox}
Düşük korelasyonu otomatik olarak ``değişkenler ilişkisiz'' diye yorumlamayın.
Ranj daralması, düşük güvenirlik, doğrusal olmayan biçim ve alt grup karışımı
katsayıyı küçültebilir. Önce veri ve ölçüm koşullarını inceleyin.
\end{mistakebox}

\section{Korelasyon ve nedensellik}
\index{nedensellik}
\index{ters nedensellik}
\index{karıştırıcı değişken}

Korelasyon iki değişkenin birlikte değiştiğini gösterir; tek başına birinin
diğerine neden olduğunu göstermez. Uyku ile kaygı arasındaki negatif ilişki
en az dört biçimde açıklanabilir:

\begin{enumerate}
  \item az uyku kaygıyı artırabilir,
  \item yüksek kaygı uykuyu azaltabilir,
  \item sınav yükü veya ekonomik stres her ikisini etkileyebilir,
  \item örnekleme, yanıtlamama veya ölçüm süreci yapay ilişki üretebilir.
\end{enumerate}

\begin{figure}[htbp]
\centering
\begin{tikzpicture}[alt={Sınav yükünün hem uyku süresini hem akademik kaygıyı etkilediği ve uyku ile kaygının ilişkili olduğu nedensel şema},
  box/.style={draw=PrismBlue,fill=PrismLight,rounded corners,align=center,minimum width=2.5cm,minimum height=0.8cm}
]
\node[box] (stress) at (0,0) {Sınav yükü};
\node[box] (sleep) at (4,0) {Uyku süresi};
\node[box] (anxiety) at (8,0) {Kaygı};
\draw[->,thick,PrismBlue] (stress)--(sleep);
\draw[->,thick,PrismBlue] (stress) to[bend right=28] (anxiety);
\draw[<->,thick,PrismSubheadingBlue] (sleep)--(anxiety);
\end{tikzpicture}
\caption{Ortak neden ve belirsiz yön içeren olası açıklama şeması.}
\end{figure}

Nedensel yorum için zaman sırası, uygun karşılaştırma, karıştırıcıların
denetlenmesi, ölçüm kalitesi ve mümkünse rastgele atama gibi ek tasarım
kanıtları gerekir. İstatistiksel kontrol bütün karıştırıcıları otomatik olarak
ortadan kaldırmaz.

\section{Kısmi korelasyon ve kontrol yanılgısı}
\index{kısmi korelasyon}
\index{kontrol değişkeni}

Kısmi korelasyon, bir veya daha fazla değişkenle doğrusal olarak ortaklaşan
kısım çıkarıldıktan sonra iki değişken arasındaki ilişkiyi özetler. Örneğin
uyku ile kaygı ilişkisi sınıf düzeyi sabit tutulduğunda incelenebilir.

Ancak bir değişkeni ``kontrol etmek'' nedensel sorunu otomatik çözmez. Araya
giren değişkeni, ortak etkinin sonucunu veya ölçüm hatalı bir kovaryatı modele
eklemek yeni yanlılık oluşturabilir. Kontrol değişkenleri yalnız $p$-değerine
veya korelasyon tablosuna göre değil, kuramsal nedensel gerekçeyle seçilmelidir.

Bu konu Bölüm 11'de doğrusal regresyon çerçevesinde ayrıntılandırılacaktır.

\section{Çok sayıda korelasyon ve seçici raporlama}
\index{çoklu test}
\index{seçici raporlama}

On değişkenin bütün ikilileri için 45 korelasyon hesaplanır. Yalnız düşük
$p$-değerlerini seçmek yanlış pozitif ve abartılı etki riski doğurur.
Araştırmacı birincil ilişkileri önceden belirlemeli, korelasyon ailesinin
büyüklüğünü açıklamalı ve gerektiğinde Holm ya da yanlış keşif oranı gibi
düzeltmeler kullanmalıdır.

Korelasyon matrisi sunulduğunda her hücrenin örneklem hacmi aynı olmayabilir.
Çift bazında eksik veri silme, farklı korelasyonların farklı öğrenci
altkümelerine dayanmasına yol açabilir. Eksik veri yöntemi açıkça
belirtilmelidir.

\section{Pearson korelasyonu için koşullar}

Korelasyon ve sıfır korelasyon testi değerlendirilirken şu noktalar incelenir:

\begin{itemize}
  \item her gözlem çifti aynı analiz birimine aittir,
  \item gözlem çiftleri birbirinden bağımsızdır veya bağımlılık modellenmiştir,
  \item iki değişken niceldir ve ilişki yaklaşık doğrusaldır,
  \item güçlü aykırı ve yüksek etkili gözlem yoktur,
  \item klasik küçük örneklem testi için ortak dağılım yaklaşık iki değişkenli normaldir,
  \item örnekleme ve ölçüm süreci hedeflenen yoruma uygundur.
\end{itemize}

Korelasyon katsayısının hesaplanması için normallik gerekmez; fakat klasik
$p$-değeri ve güven aralığının küçük örneklem geçerliği dağılım koşullarından
etkilenebilir. Grafik, bootstrap aralığı ve sağlamlık analizi bu nedenle
yararlıdır.

\begin{assumptionbox}
İki değişken için ayrı ayrı normal histogram görmek ilişkinin doğrusal veya
aykırı değersiz olduğunu göstermez. Varsayım incelemesinin temel grafiği iki
değişkeni birlikte gösteren saçılım grafiğidir.
\end{assumptionbox}

\section{Python ile yeniden üretilebilir analiz}

\begin{lstlisting}
import numpy as np
import matplotlib.pyplot as plt
from scipy import stats

uyku = np.array([4.5,5.0,5.2,5.5,5.8,6.0,6.1,6.3,6.5,6.7,
                 6.8,7.0,7.1,7.3,7.5,7.6,7.8,8.0,8.2,8.5])
kaygi = np.array([68,52,75,60,66,48,70,55,64,58,
                  45,62,53,46,60,43,55,49,57,40])

r, p = stats.pearsonr(uyku, kaygi)
rs, ps = stats.spearmanr(uyku, kaygi)

# Fisher donusumuyle yuzde 95 guven araligi
z = np.arctanh(r)
se = 1 / np.sqrt(len(uyku) - 3)
alt, ust = np.tanh([z - 1.96 * se, z + 1.96 * se])

print("Pearson r ve p:", r, p)
print("Yuzde 95 GA:", alt, ust)
print("Spearman rs ve p:", rs, ps)
print("r-kare:", r**2)

plt.scatter(uyku, kaygi)
egim, kesme = np.polyfit(uyku, kaygi, 1)
plt.plot(uyku, kesme + egim * uyku, color="navy")
plt.xlabel("Gecelik uyku suresi")
plt.ylabel("Akademik kaygi puani")
plt.show()
\end{lstlisting}

Kod çıktısından önce grafiğin biçimi ve etkili gözlemler incelenmelidir.
Korelasyon katsayıları birbirinin yerine seçilecek iki farklı ``anlamlılık
yolu'' değildir; Pearson doğrusal, Spearman sıralara dayalı tekdüze ilişkiyi
hedefler.

\section{SPSS ile korelasyon analizi}

\begin{enumerate}
  \item Önce \textbf{Graphs $>$ Chart Builder} ile saçılım grafiği oluşturun.
  \item \textbf{Analyze $>$ Correlate $>$ Bivariate} yolunu izleyin.
  \item İki değişkeni analiz alanına taşıyın; araştırma hedefine göre Pearson veya Spearman seçin.
  \item Tek veya çift yönlü seçimin veri görülmeden önceki hipotezle uyumunu doğrulayın.
  \item Eksik değerlerin çift bazında mı liste bazında mı dışlandığını kontrol edin.
  \item Katsayı, kesin $p$-değeri ve her çift için geçerli $n$ değerini okuyun.
  \item Güven aralığı sürümde doğrudan verilmiyorsa bootstrap seçeneğini veya uygun ek yordamı kullanın.
\end{enumerate}

\begin{outputguidebox}
Korelasyon tablosunu şu sırayla okuyun: (1) değişkenlerin yönü ve kodlaması,
(2) her çift için geçerli $n$, (3) katsayının işareti ve büyüklüğü, (4) kesin
$p$-değeri, (5) güven aralığı, (6) saçılım grafiği ve etkili gözlemler.
Yıldız işaretlerinin sayısını ilişkinin büyüklüğü sanmayın; yıldızlar yalnız
seçilmiş anlamlılık eşiklerini gösterir.
\end{outputguidebox}

\section{Bilimsel raporlama örneği}

``Gecelik uyku süresi ile akademik kaygı puanı arasında negatif doğrusal
ilişki bulunmuştur, $r(18)=-0{,}585$, yüzde 95 GA
$[-0{,}816;-0{,}192]$, $p=0{,}007$, $n=20$. Saçılım grafiği belirgin bir
eğrilik göstermemiştir. Spearman duyarlılık analizi benzer yönde sonuç
vermiştir, $r_s=-0{,}549$, $p=0{,}012$. Kesitsel ve öz-bildirime dayalı
tasarım nedeniyle ilişki nedensel etki olarak yorumlanmamıştır.''

Serbestlik derecesi Pearson korelasyon testinde $n-2$ olduğu için parantezde
18 verilmiştir. Alan yazındaki raporlama geleneğine göre doğrudan $n=20$ de
belirtilebilir. Yöntem, katsayı türü ve eksik veri işlemi açık olmalıdır.

\section{Bir korelasyon araştırmasını değerlendirirken}

\begin{enumerate}
  \item Araştırma sorusu doğrusal mı, tekdüze mi, yoksa daha genel bir ilişkiyi mi hedefliyor?
  \item Saçılım grafiği yön, biçim, alt gruplar ve etkili gözlemler bakımından incelenmiş mi?
  \item Pearson veya Spearman seçimi veri yapısıyla gerekçelendirilmiş mi?
  \item Katsayı güven aralığı, $p$-değeri ve örneklem hacmiyle birlikte verilmiş mi?
  \item Ranj daralması, ölçüm güvenirliği ve eksik veri dikkate alınmış mı?
  \item Çok sayıda korelasyon için seçici raporlama veya çoklu test sorunu var mı?
  \item Korelasyon açıklanan yüzde, tahmin doğruluğu veya nedensel etki gibi yanlış yorumlanmış mı?
\end{enumerate}

\section{R ile korelasyon uygulaması}
\label{sec:r-correlation}

\textbf{Soru.} Bölümdeki 20 öğrencinin uyku süresi ve kaygı puanları için
saçılım grafiğini, Pearson ve Spearman korelasyonlarını ve Pearson güven
aralığını üretin.

\begin{lstlisting}[style=prismR]
uyku <- c(4.5,5.0,5.2,5.5,5.8,6.0,6.1,6.3,6.5,6.7,
          6.8,7.0,7.1,7.3,7.5,7.6,7.8,8.0,8.2,8.5)
kaygi <- c(68,52,75,60,66,48,70,55,64,58,
           45,62,53,46,60,43,55,49,57,40)

plot(uyku, kaygi, pch = 19,
     xlab = "Uyku suresi", ylab = "Kaygi puani")
abline(lm(kaygi ~ uyku), col = "blue", lwd = 2)

pearson <- cor.test(uyku, kaygi, method = "pearson")
spearman <- cor.test(uyku, kaygi, method = "spearman",
                     exact = FALSE)
pearson
spearman

# Tek gozlemin duyarlilik etkisi
cor(uyku[-3], kaygi[-3], method = "pearson")
\end{lstlisting}

\begin{outputguidebox}
\textbf{Çözüm ve kontrol değerleri.} Pearson korelasyonu
$r=-0{,}5849$, $t(18)=-3{,}061$, $p=0{,}00675$'tir. Spearman katsayısı
yaklaşık $\rho_s=-0{,}5493$, $p=0{,}0121$ olur. Kesin güven aralığı
\texttt{cor.test} çıktısından okunur. Katsayıların yönü aynıdır; sayısal
büyüklüklerinin farklı olması bir yazılım hatası değildir.
\end{outputguidebox}

\textbf{Yorum.} Pearson doğrusal birlikte değişimi, Spearman sıraların tekdüze
ilişkisini özetler. İki katsayıyı sonuçtan sonra daha küçük $p$-değerini seçmek
için değil, veri yapısı ve araştırma hedefi doğrultusunda kullanın. Gözlemsel
kesitsel veri nedensel uyku etkisini göstermez.

\textbf{Pekiştirme ve yanıt.} Üçüncü gözlemi çıkarmak yalnız duyarlılık
karşılaştırmasıdır. Kayıt hatası veya önceden belirlenmiş bir dışlama gerekçesi
yoksa gözlem sırf katsayı değiştiği için silinmemelidir.
\section{Bölüm sonu çözümlü problemler}

\begin{enumerate}
  \item \textbf{Korelasyon testi.} $r=0{,}50$ ve $n=27$ için sıfır korelasyon testinin $t$ değerini bulun.

  \textbf{Çözüm:} $t=r\sqrt{(n-2)/(1-r^2)}=0{,}50\sqrt{25/0{,}75}\approx2{,}89$'dur. Serbestlik derecesi $n-2=25$'tir.

  \item \textbf{Açıklanan değişkenlik.} $r=-0{,}60$ sonucunu yön ve $r^2$ bakımından yorumlayın.

  \textbf{Çözüm:} İlişki negatiftir; bir değişken arttıkça diğeri doğrusal olarak azalma eğilimindedir. $r^2=0{,}36$ olup basit doğrusal modelde örneklem değişkenliğinin yüzde 36'sı ortak doğrusal değişkenliktir. Bu oran nedensel etki değildir.

  \item \textbf{Yöntem seçimi.} Saçılım grafiğinde ilişki sürekli azalan fakat belirgin eğrisel ve iki uç gözleme duyarlı görünmektedir. Pearson mı Spearman mı önceliklidir?

  \textbf{Çözüm:} Araştırma hedefi tekdüze sıralı ilişkiyse Spearman daha uygundur; doğrusal biçim gerektirmez ve uç büyüklüklerin etkisini azaltır. Yine de grafik ve etkili gözlem incelemesi raporlanmalı, Spearman otomatik bir aykırı değer çözümü sayılmamalıdır.
\end{enumerate}

\bolumkaynaklari{\cite{anscombe1973,kutner2005,james2021,cohen1988}}

\section*{Hatırla--Uygula--Yansıt}

\begin{enumerate}
\renewcommand{\labelenumi}{\textbf{\Alph{enumi}.}}
  \item \textbf{Hatırla:} Saçılım grafiğinde ilişkinin yön, biçim, güç ve aykırı değer özelliklerini açıklayın.
  \item \textbf{Yorumla:} $r=-0{,}60$ için yönü, büyüklüğü ve basit doğrusal modeldeki $R^2$ değerini belirtin.
  \item \textbf{Ayır:} Korelasyon katsayısı ile regresyon eğiminin neden aynı şey olmadığını açıklayın.
  \item \textbf{Test et:} $r=0{,}40$ ve $n=27$ için sıfır korelasyon testinin $t$ istatistiğini hesaplayın.
  \item \textbf{Aralık:} $r=0{,}35$ ve yüzde 95 GA $[-0{,}05;0{,}65]$ sonucunu kanıt ve hassasiyet bakımından yorumlayın.
  \item \textbf{Seç:} Pearson yerine Spearman kullanmayı gerektirebilecek üç veri özelliği yazın.
  \item \textbf{Eğrilik:} Güçlü U biçimli ilişkide Pearson ve Spearman katsayılarının neden sıfıra yakın olabileceğini açıklayın.
  \item \textbf{Aykırı değer:} Tek yüksek kaldıraçlı gözlemin korelasyonu nasıl büyütebileceğini veya küçültebileceğini tartışın.
  \item \textbf{Ranj:} Yalnız yüksek başarı gösteren öğrencileri seçmenin kaygı--başarı korelasyonuna olası etkisini açıklayın.
  \item \textbf{Alt grup:} Toplam korelasyon ile grup içi korelasyonların neden farklı olabileceğine bir PDR örneği verin.
  \item \textbf{Ölçüm:} Düşük güvenirliğin gözlenen korelasyonu nasıl etkileyebileceğini açıklayın.
  \item \textbf{Nedensellik:} Uyku ile kaygı arasındaki ilişki için ters yön ve iki ortak neden önerin.
  \item \textbf{Çoklu test:} Sekiz değişkenin bütün ikilileri için kaç korelasyon hesaplanacağını bulun ve seçici raporlama riskini tartışın.
  \item \textbf{Uygula:} Python verisine $(9,90)$ gözlemini ekleyip Pearson, Spearman ve grafikteki değişimi karşılaştırın.
  \item \textbf{Raporla:} Katsayı türü, $r$, serbestlik derecesi veya $n$, güven aralığı, $p$-değeri, grafik ve nedensellik sınırı içeren kısa bir bulgu paragrafı yazın.
\end{enumerate}